\section{Methods}
\label{sec:methods}

\subsection{Dataset and task}

We evaluate remaining-useful-life regression on N-CMAPSS~\cite{ariaschao2021ncmapss}, using the train, validation, and test splits exposed by the \texttt{rul\_datasets} reader~\cite{ruldatasets2026}. This is a simulation benchmark: N-CMAPSS combines a commercial-engine simulation model with real recorded flight profiles, and our experiments should not be read as validation on in-service engine failures. Each engine unit is represented as a list of per-cycle multivariate sensor sequences, and each cycle has one scalar RUL target. Unless otherwise noted, the reader caps RUL labels at 65 cycles and uses its default split, where the last 20\% of the original training units are used for validation. For FD1 this yields 5 development units, 1 validation unit, and 4 test units. The default input has shape $(T,32)$, where $T$ is the retained temporal length and the 32 channels are 4 operating-condition variables, 14 physical sensors, and 14 virtual sensors. FD1 is the primary controlled benchmark and FD2 is used as a cross-subset robustness check.

\subsection{Variable-length sequence preparation}

Raw N-CMAPSS cycles can contain long, variable-length flight records. The training pipeline flattens the per-engine lists into per-cycle examples and, when a maximum sequence length is specified, keeps the last $T$ timesteps of each cycle. Unless otherwise noted, the apples-to-apples benchmark uses $T=1000$. This last-window convention encodes the hypothesis that recent degradation behavior is more predictive of near-term RUL than distant early-flight history. After truncation, feature normalization is fit on the training split only and applied unchanged to validation and test examples, avoiding leakage from held-out engines.

\begin{figure}[t]
\centering
\includegraphics[width=\linewidth]{pipeline_overview}
\caption{Preprocessing and evaluation pipeline. Each engine is flattened into per-cycle examples, truncated to the last $N$ timesteps, normalized with train-fit feature statistics, passed through the shared model registry, and evaluated with the same RUL metrics.}
\label{fig:pipeline}
\end{figure}

\subsection{Loss and metrics}

All controlled benchmark rows in the main comparison use asymmetric MSE with late-prediction multiplier $\alpha=2.0$:
\begin{equation}
\mathcal{L} =
\frac{1}{N}\sum_{i=1}^{N}
\begin{cases}
\alpha(\hat{y}_i-y_i)^2, & \hat{y}_i-y_i \ge 0,\\
(\hat{y}_i-y_i)^2, & \hat{y}_i-y_i < 0.
\end{cases}
\end{equation}
Here over-predicting RUL is treated as a late maintenance decision and is penalized more strongly than conservative under-prediction. This is a design choice motivated by the asymmetric PHM scoring function of Saxena et al.~\cite{saxena2008damage} and recent asymmetric-loss work~\cite{mpm2024loss}. Section~\ref{sec:ablations} compares asymmetric MSE with symmetric MSE and asymmetric Huber under the same WaveNet FD1 setting.

Evaluation reports RMSE, $R^2$, PHM score, and threshold accuracies. Accuracy@20 is the percentage of predictions within $\pm20$ cycles of the true RUL. The PHM score follows the standard asymmetric exponential form, with late errors penalized more sharply than early errors. We compare our own models only under matched preprocessing and training settings; externally published numbers are used as context for the N-CMAPSS result band, not as direct head-to-head baselines.

\subsection{Model registry and benchmark harness}

The repository currently registers 22 neural architectures under a common registry interface. The catalog spans a no-temporal MLP baseline, recurrent models~\cite{hochreiter1997lstm,cho2014gru}, dilated convolutional models~\cite{bai2018tcn,oord2016wavenet}, CNN--RNN hybrids, attention models~\cite{vaswani2017attention}, transformer variants~\cite{fan2024star}, and recent multi-scale TCN/fusion architectures~\cite{ttsnet2025,xu2024msattn}. This shared interface lets the same preprocessing, optimizer, callbacks, loss, and metric code drive every architecture in the benchmark. Standard hyperparameters are units=64, dense units=32, dropout=0.2, Adam~\cite{kingma2015adam} with learning rate 0.001, 30 epochs for the controlled runs, and early stopping on validation loss.

\subsection{MSTCN architecture}

The flagship model is a multi-scale temporal convolutional network with global fusion attention (MSTCN), following the architectural family of attention-based multi-scale TCNs~\cite{xu2024msattn}. Given an input sequence $(T,32)$, MSTCN applies four parallel residual TCN branches with dilation rates $1$, $2$, $4$, and $8$. Each branch captures a different temporal receptive field over the same sensor stream. The branch outputs are then passed through a global fusion attention block that combines channel attention, temporal attention, cross-scale concatenation, and a learned fusion gate. The fused representation is pooled over time and passed through a dense regression head to produce one scalar RUL estimate.

\begin{figure}[t]
\centering
\includegraphics[width=\linewidth]{mstcn_architecture}
\caption{MSTCN architecture used in the benchmark. Parallel residual TCN branches cover four dilation scales before global fusion attention combines channel, temporal, and cross-scale information for the final RUL regression head.}
\label{fig:mstcn-architecture}
\end{figure}

For the default configuration, MSTCN uses 64 filters per branch and a 32-unit dense head, yielding roughly 150k parameters. This is larger than some baselines, so our fixed-protocol claim is limited to ``competitive'' rather than ``superior at matched parameter count.'' An iso-parameter comparison has not yet been run.
